%Version 3.1 December 2024
% See section 11 of the User Manual for version history
%
%%%%%%%%%%%%%%%%%%%%%%%%%%%%%%%%%%%%%%%%%%%%%%%%%%%%%%%%%%%%%%%%%%%%%%
%%                                                                 %%
%% Please do not use \input{...} to include other tex files.       %%
%% Submit your LaTeX manuscript as one .tex document.              %%
%%                                                                 %%
%% All additional figures and files should be attached             %%
%% separately and not embedded in the \TeX\ document itself.       %%
%%                                                                 %%
%%%%%%%%%%%%%%%%%%%%%%%%%%%%%%%%%%%%%%%%%%%%%%%%%%%%%%%%%%%%%%%%%%%%%

%%\documentclass[referee,sn-basic]{sn-jnl}% referee option is meant for double line spacing

%%=======================================================%%
%% to print line numbers in the margin use lineno option %%
%%=======================================================%%

%%\documentclass[lineno,pdflatex,sn-basic]{sn-jnl}% Basic Springer Nature Reference Style/Chemistry Reference Style

%%=========================================================================================%%
%% the documentclass is set to pdflatex as default. You can delete it if not appropriate.  %%
%%=========================================================================================%%

%%\documentclass[sn-basic]{sn-jnl}% Basic Springer Nature Reference Style/Chemistry Reference Style

%%Note: the following reference styles support Namedate and Numbered referencing. By default the style follows the most common style. To switch between the options you can add or remove “Numbered” in the optional parenthesis. 
%%The option is available for: sn-basic.bst, sn-chicago.bst%  
 
%%\documentclass[pdflatex,sn-nature]{sn-jnl}% Style for submissions to Nature Portfolio journals
%%\documentclass[pdflatex,sn-basic]{sn-jnl}% Basic Springer Nature Reference Style/Chemistry Reference Style
\documentclass[pdflatex,sn-mathphys-num]{sn-jnl}% Math and Physical Sciences Numbered Reference Style
%%\documentclass[pdflatex,sn-mathphys-ay]{sn-jnl}% Math and Physical Sciences Author Year Reference Style
%%\documentclass[pdflatex,sn-aps]{sn-jnl}% American Physical Society (APS) Reference Style
%%\documentclass[pdflatex,sn-vancouver-num]{sn-jnl}% Vancouver Numbered Reference Style
%%\documentclass[pdflatex,sn-vancouver-ay]{sn-jnl}% Vancouver Author Year Reference Style
%%\documentclass[pdflatex,sn-apa]{sn-jnl}% APA Reference Style
%%\documentclass[pdflatex,sn-chicago]{sn-jnl}% Chicago-based Humanities Reference Style

%%%% Standard Packages
%%<additional latex packages if required can be included here>

\usepackage{graphicx}%
\usepackage{multirow}%
\usepackage{amsmath,amssymb,amsfonts}%
\usepackage{amsthm}%
\usepackage{mathrsfs}%
\usepackage[title]{appendix}%
\usepackage{xcolor}%
\usepackage{textcomp}%
\usepackage{manyfoot}%
\usepackage{booktabs}%
\usepackage{listings}%
%%%%

%%%%%=============================================================================%%%%
%%%%  Remarks: This template is provided to aid authors with the preparation
%%%%  of original research articles intended for submission to journals published 
%%%%  by Springer Nature. The guidance has been prepared in partnership with 
%%%%  production teams to conform to Springer Nature technical requirements. 
%%%%  Editorial and presentation requirements differ among journal portfolios and 
%%%%  research disciplines. You may find sections in this template are irrelevant 
%%%%  to your work and are empowered to omit any such section if allowed by the 
%%%%  journal you intend to submit to. The submission guidelines and policies 
%%%%  of the journal take precedence. A detailed User Manual is available in the 
%%%%  template package for technical guidance.
%%%%%=============================================================================%%%%

%% as per the requirement new theorem styles can be included as shown below
\theoremstyle{thmstyleone}%
\newtheorem{theorem}{Theorem}%  meant for continuous numbers
%%\newtheorem{theorem}{Theorem}[section]% meant for sectionwise numbers
%% optional argument [theorem] produces theorem numbering sequence instead of independent numbers for Proposition
\newtheorem{proposition}[theorem]{Proposition}% 
%%\newtheorem{proposition}{Proposition}% to get separate numbers for theorem and proposition etc.

\theoremstyle{thmstyletwo}%
\newtheorem{example}{Example}%
\newtheorem{remark}{Remark}%

\theoremstyle{thmstylethree}%
\newtheorem{definition}{Definition}%

\raggedbottom
%%\unnumbered% uncomment this for unnumbered level heads

% Packages needed by the transferred manuscript's existing tables and floats.
\usepackage{tabularx}
\usepackage{placeins}
\usepackage{adjustbox}

% Condensed article and Online Resource 1; see articol_jurnal_changes.md.
% Replace author placeholders and supply verified declarations before submission.
\newcommand{\manuscriptfigure}[2][]{\includegraphics[#1]{#2}}

\renewcommand{\thesection}{S\arabic{section}}
\renewcommand{\thetable}{S\arabic{table}}
\renewcommand{\thefigure}{S\arabic{figure}}
\renewcommand{\theequation}{S\arabic{equation}}
\begin{document}
\begin{center}
{\Large\bfseries Online Resource 1: Supplementary Information\par}
\medskip
{\large Are 100 Queries Enough? Small-Data Adaptation of Open LLMs for Scientific Reranking\par}
\medskip
Journal of Intelligent Information Systems
\end{center}
\noindent\textbf{Authors:} Andrei-Răzvan Joldea; Cosmin Cernăzanu-Glăvan.\\
\textbf{Affiliation (both authors):} Department of Computer and Information Technology,
Politehnica University of Timișoara, Timișoara, Romania.\\
\textbf{Corresponding author:} Cosmin Cernăzanu-Glăvan.\\
\textbf{Corresponding email:} \href{mailto:cosmin.cernazanu@cs.upt.ro}{cosmin.cernazanu@cs.upt.ro}.

This supplement preserves detailed methods and supporting experiments from the full manuscript. Results retain their original numerical values. References to the main article identify its tables and sections explicitly; all other references are internal to this supplement.

\FloatBarrier
\section{Implementation and reproducibility record}
\label{supp:implementation}

We first construct an in-memory mapping from each DORIS-MAE
\texttt{abstract\_id} to a concatenated title-abstract string with the format
\texttt{"Title: {title} Abstract: {abstract}"}. This representation follows the
corpus schema and allows us to format passages efficiently.

We structure every training instance as a chat-style conversation. The user
instruction asks the model to reorder $K$ passages, indexed $[1], [2], \dots, [K]$,
by decreasing relevance to the query. The assistant response is the target
permutation, for example \texttt{[2] > [5] > \dots > [K]}. We use the following
RankGPT-style user message~\cite{Sun}:

\begin{quote}\small\itshape
You are an expert academic paper reranker. Your task is to re-order the given list of
passages (from [1] to [K]) based on their relevance to the query. Respond with only
the ranking and nothing else.
\par\vspace{2pt}
Example output for 8 passages:\newline[2] \textgreater\ [5] \textgreater\ [4] \textgreater\ [8]
\textgreater\ [6] \textgreater\ [1] \textgreater\ [3] \textgreater\ [7].
\par\vspace{2pt}
Query: {query}\newline
Passages:\newline[1] Title: {title} Abstract: {abstract}\newline
[2] Title: {title} Abstract: {abstract} \dots
\par\vspace{2pt}
Your ranking (most to least relevant):
\end{quote}

This identifier-based ranking approach follows the prompt format used by the
LitSearch GPT-4o reranker~\cite{Ajith}. We apply \emph{train-on-responses-only}
masking so that the model does not learn from prompt-specific tokens. Consequently,
the loss is computed only over tokens in the assistant turn. The selected chat
template determines the instruction and response delimiters; Gemma 3, for
example, uses \texttt{<start\_of\_turn>user} and
\texttt{<start\_of\_turn>model}.

We distinguish between the sampled \emph{training} list length,
$K_{\mathrm{train}}$ and the \emph{inference} reranking depth,
$K_{\mathrm{eval}}$. Our standard training configuration uses
$K_{\mathrm{train}}=50$, while the main LitSearch experiment reranks the top
$K_{\mathrm{eval}}=100$ first-stage candidates. At inference, we place all 100
candidates in one RankGPT-style prompt and re-index them from $[1]$ to $[100]$.
We retain the title-abstract representation used in training and truncate each
candidate to approximately 300 words. The model returns a raw permutation
of identifiers, which the parser described below validates and completes.

This setup deliberately evaluates the standard model on a list length it did not
encounter during fine-tuning. We do not use a sliding window for the main
generative rerankers. In Section~\ref{sec:list-length-mismatch}, we separate the
effect of changing $K_{\mathrm{train}}$ from the effect of changing
$K_{\mathrm{eval}}$.

To make sampling reproducible, we derive deterministic seeds for every randomized
pool-construction step from the row index and query text. When an aspect-level
instance contains more than 100 zero-scored candidates, we sample negatives from
the first 100. This choice keeps the pool connected to the higher-ranked portion
of the original DORIS-MAE candidates while retaining negative diversity. We
select hyperparameters on a temporary split of the DORIS-MAE-derived data. After
selection, we train the final models on all 1,498 filtered aspect-level instances;
LitSearch and CSFCube remain held-out evaluation benchmarks throughout.

\subsection{Output parsing and validity correction}
\label{supp:output-parsing}

The reranker must emit a permutation as bracketed candidate identifiers. Before
evaluation, we therefore apply deterministic post-processing. We extract only
identifiers of the form \texttt{[i]}, discard those outside the current candidate
window and remove duplicates while preserving their first occurrence. If the
model omits a valid candidate, we append it in its original first-stage order. We
retain the complete original order when reranking fails or the output contains no
usable identifiers. This procedure guarantees that every evaluated output is a
permutation of its input list. The metrics therefore describe the model output
together with the repair rule. Omitted candidates retain their first-stage order
at the end of the list.

\subsection{Base models fine-tuning}

For the primary cohort, we use LoRA~\cite{Hu} for parameter-efficient fine-tuning
(PEFT). The weights of Llama 3.1 8B, Gemma 3 12B and Qwen 2.5 14B are loaded
in 16-bit precision. For Gemma 3 27B, we instead use 4-bit
QLoRA~\cite{Dettmers} to control the memory
footprint. LoRA is applied to every linear projection in each language-model layer,
with rank $r = 128$, $\alpha = 256$ and dropout $0.05$. Because our task is
purely textual, we keep the vision encoders of multimodal models frozen. The LoRA
bias is disabled and the target modules are \texttt{q\_proj}, \texttt{k\_proj},
\texttt{v\_proj}, \texttt{o\_proj}, \texttt{gate\_proj}, \texttt{up\_proj} and
\texttt{down\_proj}.

We train in \texttt{bfloat16} with the Unsloth~\cite{unsloth}
\texttt{FastModel} framework and the Hugging Face TRL \texttt{SFTTrainer}. All
four primary-cohort models use AdamW, a cosine learning-rate schedule,
20 warmup steps, weight
decay $0.01$, learning rate $1 \times 10^{-5}$ and two epochs. The context length
is capped at 16,384 tokens. Loss is applied only to assistant-response tokens.
the main article, Table~2 reports the settings that differ by model.
The separate parent-grouped selection study in
the main article, Section~5.3 uses a newly selected learning rate for Llama 3.1
8B. It does not change the settings of the previous checkpoints.

For the primary-cohort Gemma 3 27B model, \texttt{load\_in\_4bit=True} invokes the pinned Unsloth and
BitsAndBytes configuration: NormalFloat4 (NF4) weights, double quantization and
\texttt{bfloat16} computation. The other three models set
\texttt{load\_in\_16bit=True} and disable 4-bit and 8-bit loading. The trained
adapters are merged into 16-bit checkpoints for local inference.

To test whether the mixed Gemma 3 27B nDCG behavior is caused by 4-bit
adaptation, we repeat its training with the corresponding
\texttt{google/gemma-3-27b-it} checkpoint, set
\texttt{load\_in\_4bit=False}, \texttt{load\_in\_8bit=False} and
\texttt{load\_in\_16bit=True}, and retain LoRA rather than full-model
fine-tuning (\texttt{full\_finetuning=False}). The control uses the same
DORIS-MAE supervision, LoRA configuration, optimization schedule, two epochs,
effective batch size and seed 3407 as the 4-bit run. We evaluate both adapters
on identical LitSearch candidate pools and input orders. The 4-bit run remains
the primary-cohort 27B configuration. Reconstructed-environment Gemma 3 27B
runs use 16-bit LoRA; the separate 16-bit run here is used as a precision control.

\subsection{Compact supervised reranker}
\label{sec:compact-training}

To assess whether a compact supervised model could offer a competitive alternative,
we fine-tune GTE-ModernBERT\cite{gte-modernbert}, a cross-encoder with
149,605,633 parameters. We initialize it from
\texttt{Alibaba-NLP/gte-reranker-modernbert-base}, pinned to revision
\texttt{204bdd4}. This checkpoint already has ranking supervision. We therefore
evaluate both the released checkpoint and its DORIS-MAE-adapted version. Here,
``zero-shot'' means without additional DORIS-MAE adaptation, not without prior
ranking training.

The compact model uses 4,494 frozen lists constructed from the DORIS-MAE
aspect queries and graded relevance labels, with the nominal sampling settings
$V=3$, $K_{\mathrm{train}}=50$ and $P=5$. The artifact preserves the sampler's
realized behavior: 4,312 lists contain 50 candidates, 182 contain 51, and 1,392
contain more than five positives. The frozen artifact contains 8,477 unique
document IDs,
whereas reconstruction from the retained LLM preparation data exposes 8,480.
Thus, the experiments share the DORIS-MAE source and nominal sampling recipe,
but exact identity with every historical LLM training list is not established.
The retained preparation and training manifests record the source and
frozen-list checksums. We do not interpret this comparison as a fully matched
training-data ablation.

We update all cross-encoder parameters in \texttt{bfloat16}, without
quantization, using LambdaLoss with NDCGLoss2++ weighting at $k=20$.
The run uses seed 3407, two epochs, learning rate $2\times10^{-5}$,
warmup ratio $0.06$ and weight decay $0.01$. Each device processes one list,
with gradient accumulation over eight lists and document mini-batches of eight.
Training passages retain the 300-word prefix, and each query-document pair is
limited to 1,024 tokens. Hyperparameters are fixed in advance. Neither LitSearch
nor CSFCube is used for checkpoint selection.

At evaluation, the model scores query-document pairs and sorts the candidate
identifiers by decreasing score. On LitSearch, both versions receive the same
frozen top-100 BM25, GRIT and Qwen3 candidate lists used by the LLMs. On
CSFCube, they score every candidate in the released facet-specific pool. The
query representation is the query paper's title and full abstract, and the
candidate representation is its title and full abstract, subject to the
1,024-token pair limit. Unlike the generative facet-conditioned prompt, this
representation does not explicitly name the facet. We use prediction batches of
64 pairs and the same benchmark metric implementations as for the LLM outputs.

\subsection{Evaluation harness and retriever integration}

We build our evaluation harness on the released LitSearch codebase. It provides
standardized evaluation scripts for BM25 and dense retrievers, together with an
LLM reranker based on RankGPT-style prompting~\cite{Sun}. We extend the
framework with local LLM inference through vLLM \cite{Kwon} and support for
additional retrievers, including the Qwen3 Embedding series~\cite{Zhang}.

We evaluate on two benchmarks:

\begin{itemize}
    \item \textbf{LitSearch}~\cite{Ajith}: The benchmark contains a corpus of
    64,183 recent ML and NLP papers from the ACL Anthology and ICLR, extracted
    from S2ORC. Its 597 queries are divided into broad and specific subsets, then
    further separated into inline-citation and author-written queries. We evaluate
    end-to-end retrieval followed by reranking of the top-100 first-stage candidates
    ($K_{\mathrm{eval}}=100$).

    \item \textbf{CSFCube}~\cite{Mysore}: The benchmark provides 50 English query-facet cases from 34 unique query
    papers in computational linguistics and machine learning venues.
    Per-query candidate pools are drawn from an S2ORC computer science corpus,
    and graded relevance judgments are produced through TREC-style depth-$k$
    pooling ($k = 100$ or $250$). Because CSFCube releases these pools directly,
    we evaluate reranking alone and omit first-stage retrieval. Relevance is
    annotated separately for background, method and result, so we rerank all three
    facets.
\end{itemize}

For LitSearch, we report Recall@5 and Recall@20. For CSFCube, we report
nDCG\%@20. These choices follow the metrics used by the authors of each
benchmark.

\subsection{Reproducibility record}
\label{sec:reproducibility}

The primary and reconstructed LLM experiments use one NVIDIA H200 NVL. The GPU
has 143,771 MiB of memory according to \texttt{nvidia-smi}. A machine snapshot
taken on August 30, 2026 reports driver 580.173.02 and CUDA 13.0. The CUDA value
shown by \texttt{nvidia-smi} describes the version supported by the driver, not
the runtime used by PyTorch. Both retained training environments contain CUDA
12.8.90 runtime packages. Because the driver may have been updated between runs
and per-run driver snapshots were not retained, we do not assign the current
driver version to every earlier checkpoint.

Table~\ref{tab:software-environments} reports the main software versions from the
retained environment exports. We keep the primary and reconstructed cohorts
separate because they were produced with different environments. Complete Conda
exports are included in the accompanying reproducibility artifact.

The compact-reranker experiment uses a separate environment on one NVIDIA H200.
Its training manifest records Python 3.11.16, PyTorch 2.13.0,
Transformers 5.16.1 and Sentence Transformers 6.0.0. The GPU snapshot reports
driver 595.71.05 and 143,771 MiB of memory. Training takes 3,274.3 seconds
(54.6 minutes), with peak allocated and reserved GPU memory of 15.10 and
16.34 GiB, respectively. Complete training, evaluation and compute-benchmark
manifests are retained separately from the LLM cohort records.

\begin{table}[t]
\centering
\scriptsize
\setlength{\tabcolsep}{4pt}
\renewcommand{\arraystretch}{1.08}
\begin{tabularx}{\textwidth}{p{0.15\textwidth} X X}
\hline
Component & Primary cohort & Reconstructed revision cohort \\
\hline
Training & Python 3.12.12; PyTorch 2.9.0; Transformers 5.5.0; TRL 0.24.0; PEFT 0.18.0; Unsloth 2026.4.5; BitsAndBytes 0.48.2; Datasets 4.3.0; Accelerate 1.11.0 & Python 3.11.15; PyTorch 2.10.0; Transformers 5.5.0; TRL 0.24.0; PEFT 0.19.1; Unsloth 2026.7.3; BitsAndBytes 0.49.2; Datasets 4.3.0; Accelerate 1.14.0 \\
CUDA runtime & \multicolumn{2}{p{0.77\textwidth}}{NVIDIA CUDA runtime 12.8.90 and cuDNN 9.10.2.21.} \\
LitSearch inference & \multicolumn{2}{p{0.77\textwidth}}{Python 3.12.12; PyTorch 2.10.0; Transformers 5.5.4; vLLM 0.19.0.} \\
CSFCube metrics & \multicolumn{2}{p{0.77\textwidth}}{Python 3.10.20; NumPy 1.26.4.} \\
\hline
\end{tabularx}
\caption{Software environments used for the reported experiments. Versions are
taken from the retained environment exports rather than inferred from current
package installations.}
\label{tab:software-environments}
\end{table}

The primary model runs use trainer and LoRA seed 3407. The temporary data split
uses seed 42. The reconstructed Llama and Gemma 3 12B seed experiments use 42,
2026 and 3407. The parent-query learning curve uses the same three seeds, while
the matched ablations use seed 42. Candidate-pool sampling is deterministic: its
seed is the MD5 hash of the source-row identifier and query text. The paired
statistical procedure uses master seed 20260815.

The tie-handling experiment uses a separate matched Llama 3.1 8B cohort with
seeds 42, 2026 and 3407. Its manifests pin the base checkpoint to revision
\texttt{0e9e39f} and record PyTorch 2.10.0, Transformers 5.5.0, PEFT 0.19.1
and Unsloth 2026.7.3 for training. Inference uses vLLM 0.19.0 with
Transformers 4.57.6, rather than the inference environment in
Table~\ref{tab:software-environments}. A validated tokenizer compatibility copy
and an explicit \texttt{rope\_theta=500000.0} override preserve the tokenizer
and rotary-position configuration across this version boundary. The retained
artifacts include the frozen inputs, raw outputs, training and evaluation
manifests, and the three-seed analysis.

For the original hyperparameter search, we split the 4,494 formatted DORIS-MAE
lists into
3,594 training, 450 validation and 450 temporary test examples using seed 42.
This was an 80/10/10 split at the formatted-list level, so aspects and augmented
pools from the same parent could occur in both training and validation. We
therefore conduct an additional learning-rate selection experiment for Llama
3.1 8B using parent-query-disjoint partitions. Within each of the five DORIS-MAE
query categories, we assign 16 parents to training, two to validation and two
to an internal test set. This gives 80/10/10 parents, with all aspects and
augmented lists from a parent kept in the same partition. The internal test
parents are fixed using seed 42; training and validation assignments use split
seeds 42, 2026 and 3407. The training seed is fixed at 3407.

We compare learning rates $5\times10^{-6}$, $1\times10^{-5}$ and
$2\times10^{-5}$, with the remaining training hyperparameters fixed. Each trial
starts from the same base checkpoint. Selection minimizes validation
negative log-likelihood (NLL): we average the response-token NLL of each list
within its parent, give equal weight to validation parents and then average
across the three splits. Internal-test losses are evaluated only after this
choice. We then refit the selected configuration from the base checkpoint on
all 4,494 lists. LitSearch and CSFCube are not used for hyperparameter selection.
They remain external evaluation datasets throughout.

The grouped-selection study uses a separate training environment. Its retained
run output reports PyTorch 2.14.1+cu130, Transformers 4.57.6 and Unsloth
2026.10.1. Training uses the Hugging Face \texttt{Trainer} with
Unsloth's \texttt{FastLanguageModel}. We report this study separately from the
primary and reconstructed cohorts; it is not an exact rerun of the historical
hyperparameter search.

\FloatBarrier
\section{Metric validation and statistical procedures}
\label{supp:metrics}

In the primary and reconstructed cohorts, we compare each fine-tuned model
with its corresponding base model on
LitSearch~\cite{Ajith} and CSFCube~\cite{Mysore}. We also position our strongest
configurations against the closed-source GPT-4o results reported by the LitSearch
authors, although we do not rerun GPT-4o in our harness. Throughout this section,
``base'' denotes the off-the-shelf instruction-tuned checkpoint, such as
Llama-3.1-8B-Instruct or gemma-3-27b-it, evaluated zero-shot with the
RankGPT-style prompt described earlier. ``FT'' denotes the same checkpoint after
our listwise fine-tuning procedure. We keep Stage 1 unchanged within every locally
evaluated base-FT pair. Differences within a pair can therefore be attributed to
the reranker.

LitSearch does not provide a script for computing retrieval metrics from ranked
outputs, so we implement one. We validate it by reproducing the retrieval-only
results reported by Ajith et al.~\cite{Ajith} for BM25, GTR-T5-large and
GritLM-7B. As Table~\ref{tab:litsearch-metric-validation} shows, our recall values
match the reported results to within rounding, confirming that our computation
follows the original procedure. We use this validated script for every LitSearch
result in the paper.

For the original LLM experiments, we distinguish two cohorts. The
\emph{primary cohort} contains
the original training checkpoints and all original main, external-baseline,
CSFCube and exploratory evaluations. Because changes in software dependencies prevented exact reproduction of the original software environment during ablation experiments, we constructed and pinned a \emph{reconstructed environment}. In this second environment, we ran
the ablations, the training-seed study and the strict document-disjoint
LitSearch evaluation, always comparing against a standard-recipe control trained
under the same conditions. We do not interpret numerical differences across the
two environments as treatment effects. For transparency,
Section~\ref{app:reconstructed-replication} reports a complete replication of the
four model families in the reconstructed environment. All reconstructed-environment
models use 16-bit LoRA, including Gemma 3 27B.
The additional parent-grouped selection study is reported separately in
the main article, Section~5.3. Its refitted checkpoint is not substituted into
the historical base--FT comparisons or ablation controls.

\begin{table}[t]
  \centering
  \scriptsize
  \setlength{\tabcolsep}{3.5pt}
  \renewcommand{\arraystretch}{1.05}
  \resizebox{\textwidth}{!}{
  \begin{adjustbox}{max width=\linewidth}
  \begin{tabular}{l c c c c c c c c c}
  \hline
  Retriever & IB R@20 & IS R@5 & IS R@20 & AB R@20 & AS R@5 & AS R@20
  & AvgB R@20 & AvgS R@5 & AvgO R@20 \\
  \hline
  BM25 (paper)        & 37.4 & 38.5 & 55.8 & 48.6 & 62.6 & 73.5 & 39.9 & 50.0 & 57.9 \\
  BM25 (ours)         & 37.4 & 38.5 & 55.8 & 48.6 & 62.6 & 73.5 & 39.9 & 50.0 & 57.9 \\
  GTR-T5-large (paper)& 45.7 & 38.5 & 51.5 & 37.1 & 40.8 & 55.9 & 43.8 & 39.6 & 51.0 \\
  GTR-T5-large (ours) & 45.8 & 38.5 & 51.5 & 37.1 & 40.8 & 55.9 & 43.8 & 39.6 & 51.1 \\
  GritLM-7B (paper)   & 69.7 & 67.7 & 77.9 & 74.3 & 82.5 & 89.1 & 70.8 & 74.8 & 80.0 \\
  GritLM-7B (ours)    & 69.7 & 67.7 & 77.9 & 74.3 & 82.5 & 89.1 & 70.8 & 74.8 & 80.0 \\
  \hline
  \end{tabular}
  \end{adjustbox}
  }
  \caption{Validation of our LitSearch metric script against the retrieval-only
  results reported by Ajith et al.~\cite{Ajith}. Values are percentages. IB/IS
  denote inline broad/specific queries, while AB/AS denote author-written
  broad/specific queries.}
  \label{tab:litsearch-metric-validation}
  \end{table}

\subsection{Evaluation metrics}

\paragraph{Recall@$k$ (R@$k$)} LitSearch pairs each query with one or more
relevant paper IDs and treats relevance as binary; it does not assign graded
relevance scores to individual papers. For each query, we compute
\(
R@k = \frac{|\text{top-}k \text{ retrieved} \cap \text{relevant}|}{|\text{relevant}|}
\)
and report the macro-average over all 597 queries. R@5 measures recall at an early
cutoff. R@20 represents a more practical retrieval-augmented setting in which a
small candidate set is passed to a downstream model or user.

\paragraph{nDCG@$k$ (LitSearch)} Because LitSearch uses binary relevance, our
harness assigns a gain of $1$ to every retrieved relevant document and $0$ to all
other documents. We then apply a logarithmic discount according to rank:
\[
\mathrm{DCG@}k = \sum_{i=1}^{k}
\frac{\mathbb{I}[d_i \in \mathrm{rel}]}{\log_2(i+1)}
\]
We obtain nDCG@$k$ by normalizing DCG@$k$ with the ideal DCG, in which
relevant documents appear as early as possible. Unlike R@$k$, nDCG@$k$ rewards
their position within the top $k$, not only their presence.

\paragraph{nDCG\%@$k$ (CSFCube)} CSFCube assigns relevance scores from 0-3,
and its candidate-pool size varies across queries. We evaluate each query over its
released pool. Following \cite{Mysore}, we report nDCG\%@$k$: nDCG computed
over the top $k\%$ of the pool rather than at one fixed rank. For a pool of size
$n$, the cutoff is $\lfloor (k/100) \cdot n \rfloor$. Thus, nDCG\%@$20$ is
equivalent to nDCG at rank 20 for a pool of 100 candidates, but to nDCG at rank
30 for a pool of 150. We use the graded labels directly as DCG gains, without
binarization and apply a logarithmic position discount. Normalization uses the
ideal DCG obtained by sorting the pool's labels in descending order, which gives
values in $[0,1]$.

\paragraph{Result reporting} For LitSearch, R@20 and nDCG@20 are the primary
metrics, while R@5 is a descriptive secondary metric. The filtered
exact-text overlap analysis reports base, FT and arithmetic FT-minus-base
differences over the same 595 queries.

\paragraph{Paired statistical uncertainty} For every paired comparison, we align
the two system outputs by evaluation-query identifier and compute the
candidate-minus-reference difference for each metric. On LitSearch, we obtain
pointwise 95\% percentile intervals from 10,000 paired query-level bootstrap
replicates. For the aggregate pool-augmentation comparisons, we first average
each query over BM25, GRIT and Qwen3 and then resample the 597 query-level
averages. On CSFCube, the 50 query-facet cases originate from 34 query papers.
We therefore use a clustered bootstrap over unique papers, retaining all facets
of a sampled paper in the same replicate.

We calculate two-sided $p$-values with 10,000 paired sign-flip permutations.
Signs are applied at query level for LitSearch and at unique-paper level for
CSFCube. We use the plus-one Monte Carlo correction and master seed 20260815;
comparison-specific bootstrap and permutation seeds are derived
deterministically from it. Holm--Bonferroni correction is applied separately
within each metric to four prespecified families: the 12 LitSearch base--FT
comparisons, six local-baseline comparisons, two pool-augmentation comparisons
and four CSFCube base--FT comparisons. We report the raw value as $p$ and the
adjusted value as $p_{\mathrm{H}}$. The confidence intervals are pointwise and
are not multiplicity-adjusted.

For the compact reranker, we use four additional comparison families. On
LitSearch, the three Llama-FT-minus-adapted-GTE comparisons form one family
per metric, and the three adapted-GTE-minus-zero-shot-GTE comparisons form
another. On CSFCube, each corresponding contrast contains one pooled,
paper-clustered comparison per metric. Holm correction is applied within these
families separately from the earlier LLM comparisons. The Llama reference is the
primary-cohort FT output. These tests quantify evaluation-query uncertainty for
the retained checkpoints. They do not remove differences in pretraining,
training-list provenance, objectives or software environments.

\FloatBarrier
\section{LitSearch gains and precision control}
\label{supp:precision}

\begin{figure}[ht]
  \centering
  \begingroup
  \newcommand{\gaincell}[2]{\colorbox{#1}{\makebox[1.18cm][c]{#2}}}
  \newcommand{\gaindark}[2]{\colorbox{#1}{\makebox[1.18cm][c]{\color{white}#2}}}
  \setlength{\tabcolsep}{3pt}
  \renewcommand{\arraystretch}{1.08}
  \scriptsize
  \begin{adjustbox}{max width=\linewidth}
  \begin{tabular}{l c c c}
  \hline
  Model / retriever & R@20 & R@5 & nDCG@20 \\
  \hline
  Llama 3.1 8B / BM25  & \gaindark{red!58}{+11.7} & \gaindark{red!62}{+12.3} & \gaincell{red!46}{+9.1} \\
  Llama 3.1 8B / GRIT  & \gaincell{red!37}{+7.3} & \gaindark{red!67}{+13.3} & \gaindark{red!51}{+10.2} \\
  Llama 3.1 8B / Qwen3 & \gaincell{red!36}{+7.2}  & \gaindark{red!81}{+16.2} & \gaindark{red!64}{+12.7} \\
  Gemma 3 12B / BM25   & \gaincell{red!31}{+6.1}  & \gaincell{red!8}{+1.6}   & \gaincell{red!5}{+0.3} \\
  Gemma 3 12B / GRIT   & \gaincell{red!7}{+1.3}   & \gaincell{red!5}{+0.5}   & \gaincell{red!5}{+1.0} \\
  Gemma 3 12B / Qwen3  & \gaincell{red!18}{+3.6}  & \gaincell{red!8}{+1.6}   & \gaincell{red!18}{+3.5} \\
  Qwen 2.5 14B / BM25  & \gaincell{red!23}{+4.5}  & \gaincell{blue!5}{-0.5}  & \gaincell{red!5}{+0.9} \\
  Qwen 2.5 14B / GRIT  & \gaincell{red!16}{+3.1}  & \gaincell{red!5}{+0.8}   & \gaincell{red!16}{+3.1} \\
  Qwen 2.5 14B / Qwen3 & \gaincell{red!9}{+1.8}   & \gaincell{red!8}{+1.5}   & \gaincell{red!12}{+2.3} \\
  Gemma 3 27B / BM25   & \gaincell{red!5}{+0.4}   & \gaincell{blue!5}{-0.1}  & \gaincell{blue!11}{-2.2} \\
  Gemma 3 27B / GRIT   & \gaincell{red!6}{+1.1}   & \gaincell{blue!5}{-0.9}  & \gaincell{blue!5}{-0.7} \\
  Gemma 3 27B / Qwen3  & \gaincell{red!23}{+4.6}  & \gaincell{red!8}{+1.6}   & \gaincell{red!5}{+0.7} \\
  \hline
  \end{tabular}
  \end{adjustbox}
  \endgroup
  \caption{LitSearch fine-tuning gains. We compute each value as the difference between the fine-tuned reranker and its base instruction-tuned model. Red indicates an improvement and blue a regression; darker shading represents a larger absolute change.}
  \label{fig:litsearch-gains}
  \end{figure}

\begin{table}[t]
\centering
\scriptsize
\setlength{\tabcolsep}{7pt}
\renewcommand{\arraystretch}{1.08}
\begin{adjustbox}{max width=\linewidth}
\begin{tabular}{lrrrrrr}
\hline
Retriever & Base & 4-bit & 16-bit & $\Delta_{16-4}$ & 95\% CI & $p_{\mathrm{H}}$ \\
\hline
BM25  & 51.89 & 49.69 & 50.61 & $+0.92$ & [$-$0.08, 1.95] & .2328 \\
GRIT  & 65.67 & 65.01 & 64.74 & $-$0.26 & [$-$1.29, 0.76] & .6066 \\
Qwen3 & 65.22 & 65.90 & 65.02 & $-$0.88 & [$-$1.89, 0.10] & .2328 \\
\hline
\end{tabular}
\end{adjustbox}
\caption{Gemma 3 27B LitSearch nDCG@$20$ precision control. Base, 4-bit
QLoRA and 16-bit LoRA scores are percentages. Each delta is 16-bit LoRA minus
4-bit QLoRA over the same 597 queries. We report paired percentile 95\%
bootstrap intervals and Holm-corrected sign-flip permutation $p$-values.}
\label{tab:gemma27b-precision}
\end{table}

Higher precision does not consistently improve nDCG@$20$. Relative to 4-bit
QLoRA, 16-bit LoRA gains $0.92$ pp with BM25 but changes GRIT and Qwen3 by
$-0.26$ and $-0.88$ pp. All three intervals include zero, and none of the
comparisons is significant after Holm correction. The mixed top-rank behavior
therefore cannot be attributed to 4-bit quantization alone. This control covers
LitSearch only, so we do not use it to explain the CSFCube regression.

\FloatBarrier
\section{Numerical supervision-budget learning curve}
\label{supp:budget}

\begin{table}[ht]
\centering
\scriptsize
\setlength{\tabcolsep}{4.2pt}
\renewcommand{\arraystretch}{1.08}
\begin{adjustbox}{max width=\linewidth}
\begin{tabular}{c c c c c c}
\hline
Parent queries & Aspect rows & Optimizer steps & R@20 & R@5 & nDCG@20 \\
\hline
0 (zero-shot) & 0 & 0 & 66.33 & 48.49 & 47.25 \\
10  & $150.3\,[144,160]$   & $451.3\,[432,480]$   & $75.68\pm0.32$ & $63.54\pm1.67$ & $60.26\pm0.51$ \\
25  & $381.0\,[348,406]$   & $1143.3\,[1044,1218]$ & $76.26\pm0.25$ & $\mathbf{65.59\pm0.42}$ & $\mathbf{60.79\pm0.21}$ \\
50  & $746.7\,[705,782]$   & $2240.7\,[2116,2346]$ & $75.83\pm0.17$ & $64.79\pm0.86$ & $60.02\pm0.55$ \\
75  & $1123.0\,[1108,1137]$ & $3369.3\,[3324,3412]$ & $76.27\pm0.50$ & $64.52\pm0.48$ & $59.79\pm0.47$ \\
100 & $1498.0\,[1498,1498]$ & $4494.0\,[4494,4494]$ & $\mathbf{76.71\pm0.06}$ & $64.57\pm0.25$ & $60.09\pm0.08$ \\
\hline
\end{tabular}
\end{adjustbox}
\caption{Retriever-averaged parent-query learning curve. Training rows report
mean $\pm$ sample standard deviation over three seeds; aspect rows and optimizer
steps report mean $[\min,\max]$. The zero-shot row is a single checkpoint and is
not included in the seed statistics. Bold marks the best trained budget for each
metric. Values are percentages. Differences are calculated from unrounded
scores and rounded only for presentation; they may therefore differ slightly
from differences between displayed values.}
\label{tab:parent-budget-main}
\end{table}

\FloatBarrier
\section{Pool-augmentation uncertainty}
\label{supp:augmentation}

The configuration ablation and the example-count-matched control are reported in the main article. The following table gives the full paired uncertainty for the configuration comparisons.

\begin{table}[t]
\centering
\scriptsize
\setlength{\tabcolsep}{4pt}
\renewcommand{\arraystretch}{1.0}
\begin{adjustbox}{max width=\linewidth}
\begin{tabular}{l l r c c c}
\hline
Comparison & Metric & $\Delta$ & 95\% CI & $p$ & $p_{\mathrm{H}}$ \\
\hline
$V=3$ $-$ $V=1$ & R@5 & $-$0.79 & [$-$2.07, 0.47] & .2164 & .2470 \\
$V=3$ $-$ $V=1$ & R@20 & 0.95 & [$-$0.06, 1.95] & .0717 & .0717 \\
$V=3$ $-$ $V=1$ & nDCG@20 & 0.13 & [$-$0.69, 0.96] & .7664 & 1.0000 \\
$V=3$ $-$ no pool augmentation & R@5 & 0.99 & [$-$0.27, 2.22] & .1235 & .2470 \\
$V=3$ $-$ no pool augmentation & R@20 & 2.35 & [1.29, 3.40] & $<.001$ & $<.001$ \\
$V=3$ $-$ no pool augmentation & nDCG@20 & 0.19 & [$-$0.63, 1.03] & .6572 & 1.0000 \\
\hline
\end{tabular}
\end{adjustbox}
\caption{Paired uncertainty for the pool-augmentation comparisons. For each
query, metrics are first averaged over BM25, GRIT and Qwen3. Deltas are standard
$V=3$ minus the named reference in percentage points. Holm correction is applied
within the two comparisons for each metric.}
\label{tab:pool-augmentation-uncertainty}
\end{table}

\FloatBarrier
\section{Training-seed stability}
\label{supp:seeds}

To distinguish evaluation-query uncertainty from training stochasticity, we use
the three $B=100$ Llama 3.1 8B checkpoints from the parent-budget experiment
and three standard reconstructed-environment Gemma 3 12B runs.
Table~\ref{tab:seed-stability} reports the mean and sample standard deviation
over the resulting checkpoints. These deviations quantify training-run
variability only in the reconstructed environment. They are not uncertainty
estimates for the original checkpoints in the main article, Table~3.

\begin{table}[ht]
\centering
\footnotesize
\setlength{\tabcolsep}{3.5pt}
\renewcommand{\arraystretch}{1.08}
\begin{adjustbox}{max width=\linewidth}
\begin{tabular}{l l c c c}
\hline
Model & Retriever & R@20 & R@5 & nDCG@20 \\
\hline
\multirow{3}{*}{Llama 3.1 8B}
 & BM25  & $67.21\pm0.32$ & $55.58\pm0.44$ & $50.83\pm0.27$ \\
 & GRIT  & $81.93\pm0.18$ & $69.50\pm0.25$ & $64.77\pm0.11$ \\
 & Qwen3 & $80.98\pm0.28$ & $68.63\pm0.73$ & $64.68\pm0.14$ \\
\hline
\multirow{3}{*}{Gemma 3 12B}
 & BM25  & $67.24\pm0.63$ & $54.93\pm0.98$ & $49.79\pm0.18$ \\
 & GRIT  & $80.72\pm0.51$ & $69.72\pm0.50$ & $64.30\pm0.41$ \\
 & Qwen3 & $81.53\pm0.70$ & $69.19\pm0.51$ & $64.66\pm0.52$ \\
\hline
\end{tabular}
\end{adjustbox}
\caption{LitSearch stability of the standard reconstructed-environment recipe over
three training runs. We report the mean $\pm$ sample standard deviation in
percentage points.}
\label{tab:seed-stability}
\end{table}

The standard-recipe improvements retain the same direction across runs, but
seed variability remains material when an ablation gap is small. For Llama,
the retriever-specific nDCG@20 standard deviations range from 0.11 to 0.27
points. These values reinforce the caution around sub-point augmentation
differences. By contrast, the multi-point decline without aspect decomposition
in both the original and example-count-matched controls remains substantially
larger than the observed seed spread.

\FloatBarrier
\section{Inference-depth sensitivity}
\label{supp:depth}

\label{sec:list-length-mismatch}

In the main article, Table~6, the $K$ ablation changes
$K_{\mathrm{train}}$ while keeping $K_{\mathrm{eval}}=100$. We isolate inference
depth separately by fixing representative reconstructed-environment Llama 3.1
8B, Gemma 3 12B and Gemma 3 27B checkpoints and varying only the number of first-stage
candidates in the prompt. For every retriever, the shorter candidate lists are
exact prefixes of the corresponding 100-document list. Table~\ref{tab:inference-depth}
reports the equal-weight mean over BM25, GRIT and Qwen3; Table~\ref{tab:inference-depth-detailed}
gives the retriever-specific results.

\begin{table}[ht]
\centering
\footnotesize
\setlength{\tabcolsep}{4pt}
\renewcommand{\arraystretch}{1.08}
\begin{adjustbox}{max width=\linewidth}
\begin{tabular}{l c c c c}
\hline
Model & $K_{\mathrm{eval}}$ & R@20 & R@5 & nDCG@20 \\
\hline
\multirow{4}{*}{Llama 3.1 8B FT}
 & 10  & 67.15 & 62.25 & 56.03 \\
 & 20  & 72.82 & 62.74 & 57.93 \\
 & 50  & 76.30 & 63.94 & 59.36 \\
 & 100 & 77.34 & 64.93 & 60.21 \\
\hline
\multirow{4}{*}{Gemma 3 12B FT}
 & 10  & 67.15 & 62.89 & 55.58 \\
 & 20  & 72.82 & 63.18 & 57.51 \\
 & 50  & 75.80 & 63.89 & 58.82 \\
 & 100 & 77.08 & 64.13 & 59.38 \\
\hline
\multirow{4}{*}{Gemma 3 27B FT}
 & 10  & 67.15 & 63.77 & 56.63 \\
 & 20  & 72.82 & 64.84 & 58.72 \\
 & 50  & 76.80 & 65.68 & 60.08 \\
 & 100 & 77.25 & 66.29 & 60.12 \\
\hline
\end{tabular}
\end{adjustbox}
\caption{Representative reconstructed-environment LitSearch performance of fine-tuned checkpoints with $K_{\mathrm{train}}=50$ at four inference depths.
Values are percentages averaged equally over BM25, GRIT and Qwen3 for the same
597 queries. These runs are separate from the original-cohort main results. The
Gemma 3 27B checkpoint is the reconstructed 16-bit LoRA run with seed 2026, not
the primary-cohort 4-bit checkpoint.}
\label{tab:inference-depth}
\end{table}

At $K_{\mathrm{eval}}\in\{10,20\}$, R@20 is candidate-membership bounded: the
reranker cannot introduce a document outside the supplied list, and R@20 includes
every supplied candidate. The identical R@20 values across the three models at these
depths therefore follow from the shared first-stage lists, rather than identical
ranking quality. R@5 and nDCG@20 remain sensitive to the reranked order.

Increasing depth from 20 to 50 improves mean R@20 by 3.48, 2.98 and 3.98 pp for
Llama, Gemma 3 12B and Gemma 3 27B, while nDCG@20 improves by 1.42, 1.31 and
1.37 pp. Increasing depth from 50 to 100 produces smaller R@20 gains of 1.04,
1.28 and 0.45 pp and nDCG@20 gains of 0.85, 0.55 and 0.04 pp. For Gemma 3 27B,
nDCG@20 therefore effectively plateaus at depth 50. Performance has not plateaued
at depth 20, but depth 50 captures most of the improvement observed beyond 20. We
did not measure
$K_{\mathrm{eval}}=200$; a 200-document prompt would require substantially more
prefill compute, and we make no claim about that setting.

\FloatBarrier
\section{Uncertainty for local open-reranker comparisons}
\label{supp:local}

\begin{table}[t]
\centering
\scriptsize
\setlength{\tabcolsep}{4pt}
\renewcommand{\arraystretch}{0.98}
\begin{adjustbox}{max width=\linewidth}
\begin{tabular}{l l l r c c c}
\hline
Retriever & Reference & Metric & $\Delta$ & 95\% CI & $p$ & $p_{\mathrm{H}}$ \\
\hline
GRIT & Retrieval only & R@20 & 2.76 & [0.67, 4.86] & .0090 & .0360 \\
GRIT & Retrieval only & nDCG@20 & 2.95 & [1.61, 4.32] & $<.001$ & $<.001$ \\
GRIT & RankZephyr & R@20 & 1.68 & [$-$0.33, 3.78] & .1120 & .1120 \\
GRIT & RankZephyr & nDCG@20 & 0.56 & [$-$1.24, 2.38] & .5410 & .8063 \\
GRIT & Qwen3-Reranker-8B & R@20 & 2.54 & [0.11, 4.97] & .0491 & .0982 \\
GRIT & Qwen3-Reranker-8B & nDCG@20 & 11.00 & [8.32, 13.67] & $<.001$ & $<.001$ \\
\hline
Qwen3 & Retrieval only & R@20 & 3.55 & [1.68, 5.47] & $<.001$ & .0025 \\
Qwen3 & Retrieval only & nDCG@20 & 3.01 & [1.72, 4.35] & $<.001$ & $<.001$ \\
Qwen3 & RankZephyr & R@20 & 2.53 & [0.63, 4.44] & .0113 & .0360 \\
Qwen3 & RankZephyr & nDCG@20 & 0.77 & [$-$1.00, 2.62] & .4032 & .8063 \\
Qwen3 & Qwen3-Reranker-8B & R@20 & 4.87 & [2.39, 7.41] & $<.001$ & $<.001$ \\
Qwen3 & Qwen3-Reranker-8B & nDCG@20 & 12.89 & [10.20, 15.63] & $<.001$ & $<.001$ \\
\hline
\end{tabular}
\end{adjustbox}
\caption{Paired uncertainty for the local-baseline comparisons. Deltas are Gemma
3 27B FT minus the named reference in percentage points over the same 597
LitSearch queries. The intervals are pointwise; $p_{\mathrm{H}}$ is corrected
within the six local-baseline comparisons for each metric.}
\label{tab:open-reranker-uncertainty}
\end{table}

% \FloatBarrier
% \section{External GPT-4o reference comparison}
% \label{supp:gpt4o}

% To position the fine-tuned open-weight models against a strong closed-source
% reference, we use the aggregate title-abstract results reported for GPT-4o by
% Ajith et al.~\cite{Ajith}. We did not rerun GPT-4o. Our local artifacts do not
% contain the exact candidate files, API model snapshot, prompt execution, or output
% parser used to produce the published values. The comparison therefore matches the
% reported retriever family and metrics, but it does not use a shared harness and is
% not a significance-tested reranker comparison.

% The LitSearch authors report average Broad R@$20$ and average Specific R@$5$,
% but not a combined Avg.\ Specific R@20. They do provide results for both specific
% subsets. We therefore calculate their query-weighted mean using the weighting
% underlying the LitSearch values: $n=231$ inline-specific and $n=211$
% author-specific queries. This gives 85.9\% for GPT-4o + GritLM and 74.9\% for
% GPT-4o + BM25. We compute our Specific R@20 results in the same way.

% \begin{table}[ht]
% \centering
% \renewcommand{\arraystretch}{1.1}
% \setlength{\tabcolsep}{4pt}
% \footnotesize
% \begin{adjustbox}{max width=\linewidth}
% \begin{tabular}{l c c c}
% \hline
% Configuration & Broad R@20 & Spec.\ R@5 & Spec.\ R@20 \\
% \hline
% \multicolumn{4}{l}{\textit{Closed-source reference~\cite{Ajith}}} \\
% GPT-4o reranking (BM25)        & 59.9 & 68.0 & 74.9 \\
% GPT-4o reranking (GritLM)      & 75.3 & 79.2 & 85.9 \\
% \hline
% \multicolumn{4}{l}{\textit{Matched BM25 first-stage (ours)}} \\
% Llama 3.1 8B FT (BM25)         & 53.2 & 57.9 & 70.8 \\
% Gemma 3 12B FT (BM25)          & 53.8 & 58.7 & 70.5 \\
% Qwen 2.5 14B FT (BM25)         & 48.6 & 57.4 & 70.9 \\
% Gemma 3 27B FT (BM25)          & 52.1 & 62.0 & 71.2 \\
% \hline
% \multicolumn{4}{l}{\textit{Matched GritLM first-stage (ours)}} \\
% Llama 3.1 8B FT (GRIT)     & 70.4 & 72.3 & 82.6 \\
% Gemma 3 12B FT (GRIT)      & 72.4 & 75.6 & 84.3 \\
% Qwen 2.5 14B FT (GRIT)     & 73.3 & 75.5 & 83.3 \\
% Gemma 3 27B FT (GRIT)      & 73.0 & 75.8 & \textbf{86.2} \\
% \hline
% \multicolumn{4}{l}{\textit{Strong open reranker baseline}} \\
% RankZephyr (GRIT)          & 72.0 & 74.0 & 84.3 \\
% RankZephyr (Qwen3)         & 74.6 & 74.4 & 83.9 \\
% \hline
% \multicolumn{4}{l}{\textit{Retrieval-only comparison}} \\
% Qwen3-Embedding-8B & 69.7 & 74.2 & 84.3 \\
% \hline
% \multicolumn{4}{l}{\textit{Local end-to-end pipeline with stronger first stage}} \\
% Gemma 3 27B FT (Qwen3)     & \textbf{77.6} & 76.9 & \textbf{86.3} \\
% \hline
% \end{tabular}
% \end{adjustbox}
% \caption{ Comparison with the closed-source GPT-4o reranker reported by~\cite{Ajith}. “Spec. R@20” is the query-weighted average
% over inline-specific and author-specific queries, matching the breakdown used by LitSearch. The BM25 and GritLM rows
% support matched first-stage reranker comparisons; RankZephyr is included as a strong open reranker baseline; the Qwen3 rows
% are end-to-end pipeline comparisons with a stronger first-stage retriever. Bold marks values that match or exceed the strongest
% GPT-4o configuration (GritLM first-stage).}
% \label{tab:gpt4o}

% \end{table}

% Table~\ref{tab:gpt4o} separates retriever-matched reference positioning from
% retrieval-only and end-to-end pipeline results. The controlled local comparison
% with open rerankers appears instead in
% the main article, Table~11. With BM25, GPT-4o remains ahead of all
% fine-tuned open rerankers on the three reported metrics. When we match the
% GritLM retriever family, fine-tuned Gemma 3 27B slightly exceeds the reported
% GPT-4o Specific R@$20$ ($86.2\%$ vs.\ $85.9\%$), but remains lower on Broad
% R@$20$ ($-2.3$ pp) and Specific R@$5$ ($-3.4$ pp). The open reranker is therefore
% competitive at the broader R@$20$ cutoff, but it does not outperform GPT-4o
% uniformly. We cannot compute confidence intervals or significance tests for this
% cross-protocol comparison because per-query GPT-4o outputs are unavailable.

% RankZephyr provides a second open reranking reference. On the controlled local
% GritLM candidate pool, it obtains $72.0\%$ Broad R@$20$, $74.0\%$ Specific
% R@$5$ and $84.3\%$ Specific R@$20$. These results place it close to the stronger
% fine-tuned open rerankers, but below our Gemma 3 27B FT result on Specific
% R@$20$.

% The Qwen3 rows in Table~\ref{tab:gpt4o} and
% Fig.~\ref{fig:gpt4o-comparison} require a different interpretation: they compare
% complete pipelines. Fine-tuned Gemma 3 27B with the stronger Qwen3 retriever
% achieves $77.6\%$ Broad R@$20$ and $86.3\%$ Specific R@$20$, exceeding the
% reported GPT-4o + GritLM pipeline on both R@$20$ metrics. This difference reflects
% the combined contribution of a stronger first-stage retriever and the fine-tuned
% reranker. It is not a pure comparison of reranker quality against GPT-4o.

% \begin{figure}[htbp]
%   \centering
%   \manuscriptfigure[width=\columnwidth]{fig_gpt4o_comparison.pdf}
%   \caption{Comparison between the strongest closed-source GPT4o reranking baseline reported by LitSearch, Qwen3 retrieval
% without reranking, and our best open end-to-end pipeline.}
%   \label{fig:gpt4o-comparison}
%   \end{figure}

\FloatBarrier
\section{Examples of top-rank precision failures}
\label{supp:failures}

\begin{table}[htbp]
\centering
\scriptsize
\setlength{\tabcolsep}{3.5pt}
\renewcommand{\arraystretch}{1.05}
\resizebox{\textwidth}{!}{
\begin{adjustbox}{max width=\linewidth}
\begin{tabular}{p{0.12\textwidth} p{0.44\textwidth} c c c c}
\hline
Model & Query & R@20 & R@5 & nDCG@5 & Relevant positions \\
\hline
Qwen 2.5 14B &
Can you point me to studies that explore the impact of different data augmentation strategies, such as feature/token/span cutoff or dropout, in the context of contrastive learning for sentence representations? &
0.50 $\rightarrow$ 1.00 & 0.50 $\rightarrow$ 0.50 & 0.6131 $\rightarrow$ 0.2641 & [1] $\rightarrow$ [4, 20] \\
Gemma 3 12B &
What techniques have been investigated to enhance multimodal sentiment analysis by fusing different modalities? &
0.50 $\rightarrow$ 1.00 & 0.50 $\rightarrow$ 0.50 & 0.3869 $\rightarrow$ 0.2372 & [2] $\rightarrow$ [5, 12] \\
Gemma 3 27B &
Can you suggest literature on enhanced semantic parsing methods that focus on generating high-quality meaning representations and utilize knowledge-constrained decoding under specific grammar rules? &
0.3333 $\rightarrow$ 0.6667 & 0.3333 $\rightarrow$ 0.3333 & 0.2961 $\rightarrow$ 0.1815 & [2] $\rightarrow$ [5, 12] \\
\hline
\end{tabular}
\end{adjustbox}
}
\caption{LitSearch examples in which fine-tuning improves Recall@20 but
worsens the early ranking. All examples use Qwen3-Embedding-8B as the first-stage
retriever.}
\label{tab:top-rank-failures}

\end{table}

\subsection{Top-rank precision failure modes}

Although fine-tuning generally improves broader recall, it does not improve the
first positions for every query. Table~\ref{tab:top-rank-failures} presents
LitSearch examples in which the fine-tuned model retrieves more relevant papers
within the top 20, but places the earliest relevant paper lower within the top five.

The current supervised objective can therefore improve the inclusion of relevant
papers in a wider candidate window while still making consequential errors at the
first few ranks. This distinction is important in scientific search, where users
often inspect only the earliest results.

\FloatBarrier
\section{Detailed Subset and Facet Results}
\label{app:detailed-results}

\begin{table}[htbp]
\centering
\scriptsize
\setlength{\tabcolsep}{4pt}
\renewcommand{\arraystretch}{1.02}
\begin{adjustbox}{max width=\linewidth}
\begin{tabular}{l c l c c c}
\hline
Model & $K_{\mathrm{eval}}$ & Retriever & R@20 & R@5 & nDCG@20 \\
\hline
\multirow{12}{*}{Llama 3.1 8B FT}
 & \multirow{3}{*}{10}  & BM25  & 51.49 & 48.83 & 42.88 \\
 &                       & GRIT  & 74.83 & 68.91 & 62.50 \\
 &                       & Qwen3 & 75.14 & 69.02 & 62.70 \\
 & \multirow{3}{*}{20}  & BM25  & 57.94 & 50.42 & 45.57 \\
 &                       & GRIT  & 80.01 & 69.08 & 64.15 \\
 &                       & Qwen3 & 80.50 & 68.71 & 64.08 \\
 & \multirow{3}{*}{50}  & BM25  & 64.97 & 53.69 & 48.79 \\
 &                       & GRIT  & 82.78 & 70.25 & 64.96 \\
 &                       & Qwen3 & 81.14 & 67.87 & 64.31 \\
 & \multirow{3}{*}{100} & BM25  & 67.65 & 55.37 & 50.52 \\
 &                       & GRIT  & 82.44 & 70.14 & 64.99 \\
 &                       & Qwen3 & 81.93 & 69.30 & 65.11 \\
\hline
\multirow{12}{*}{Gemma 3 12B FT}
 & \multirow{3}{*}{10}  & BM25  & 51.49 & 48.41 & 42.40 \\
 &                       & GRIT  & 74.83 & 69.80 & 61.70 \\
 &                       & Qwen3 & 75.14 & 70.44 & 62.64 \\
 & \multirow{3}{*}{20}  & BM25  & 57.94 & 50.51 & 45.38 \\
 &                       & GRIT  & 80.01 & 69.69 & 63.47 \\
 &                       & Qwen3 & 80.50 & 69.35 & 63.68 \\
 & \multirow{3}{*}{50}  & BM25  & 64.81 & 53.61 & 48.54 \\
 &                       & GRIT  & 81.47 & 69.36 & 64.10 \\
 &                       & Qwen3 & 81.12 & 68.71 & 63.84 \\
 & \multirow{3}{*}{100} & BM25  & 67.91 & 53.94 & 49.70 \\
 &                       & GRIT  & 81.27 & 69.19 & 64.11 \\
 &                       & Qwen3 & 82.07 & 69.27 & 64.31 \\
\hline
\multirow{12}{*}{Gemma 3 27B FT}
 & \multirow{3}{*}{10}  & BM25  & 51.49 & 49.56 & 43.49 \\
 &                       & GRIT  & 74.83 & 71.65 & 63.02 \\
 &                       & Qwen3 & 75.14 & 70.11 & 63.37 \\
 & \multirow{3}{*}{20}  & BM25  & 57.94 & 51.68 & 46.55 \\
 &                       & GRIT  & 80.01 & 71.56 & 64.49 \\
 &                       & Qwen3 & 80.50 & 71.28 & 65.11 \\
 & \multirow{3}{*}{50}  & BM25  & 64.81 & 55.62 & 49.53 \\
 &                       & GRIT  & 82.78 & 70.56 & 65.19 \\
 &                       & Qwen3 & 82.82 & 70.86 & 65.53 \\
 & \multirow{3}{*}{100} & BM25  & 67.07 & 56.54 & 50.61 \\
 &                       & GRIT  & 82.28 & 71.31 & 64.74 \\
 &                       & Qwen3 & 82.40 & 71.03 & 65.02 \\
\hline
\end{tabular}
\end{adjustbox}
\caption{Retriever-specific inference-depth results underlying
Table~\ref{tab:inference-depth}. The fine-tuned reconstructed-environment
checkpoints are held fixed while $K_{\mathrm{eval}}$ varies. All values are
percentages over the same 597 LitSearch queries.}
\label{tab:inference-depth-detailed}
\end{table}

\begin{table}[htbp]
\centering
\scriptsize
\setlength{\tabcolsep}{4.0pt}
\renewcommand{\arraystretch}{1.04}
\begin{adjustbox}{max width=\linewidth}
\begin{tabular}{c l c c c}
\hline
Parent queries & Retriever & R@20 & R@5 & nDCG@20 \\
\hline
\multirow{3}{*}{10}
 & BM25  & $66.40\pm0.35$ & $55.56\pm0.21$ & $51.81\pm0.32$ \\
 & GRIT  & $80.07\pm0.92$ & $67.46\pm2.64$ & $64.31\pm1.04$ \\
 & Qwen3 & $80.56\pm0.59$ & $67.59\pm2.25$ & $64.67\pm0.71$ \\
\hline
\multirow{3}{*}{25}
 & BM25  & $66.86\pm0.64$ & $55.82\pm0.56$ & $51.25\pm0.23$ \\
 & GRIT  & $81.01\pm0.57$ & $70.72\pm0.12$ & $65.04\pm0.20$ \\
 & Qwen3 & $80.92\pm0.83$ & $70.25\pm0.73$ & $66.07\pm0.30$ \\
\hline
\multirow{3}{*}{50}
 & BM25  & $66.45\pm0.21$ & $56.27\pm1.49$ & $50.77\pm0.65$ \\
 & GRIT  & $80.36\pm0.68$ & $69.57\pm0.43$ & $64.10\pm0.46$ \\
 & Qwen3 & $80.67\pm0.18$ & $68.53\pm1.03$ & $65.17\pm0.60$ \\
\hline
\multirow{3}{*}{75}
 & BM25  & $66.61\pm0.20$ & $54.51\pm0.90$ & $50.06\pm0.64$ \\
 & GRIT  & $81.41\pm0.94$ & $69.76\pm0.67$ & $64.31\pm0.69$ \\
 & Qwen3 & $80.80\pm0.69$ & $69.30\pm0.08$ & $65.00\pm0.30$ \\
\hline
\multirow{3}{*}{100}
 & BM25  & $67.21\pm0.32$ & $55.58\pm0.44$ & $50.83\pm0.27$ \\
 & GRIT  & $81.93\pm0.18$ & $69.50\pm0.25$ & $64.77\pm0.11$ \\
 & Qwen3 & $80.98\pm0.28$ & $68.63\pm0.73$ & $64.68\pm0.14$ \\
\hline
\end{tabular}
\end{adjustbox}
\caption{Retriever-specific parent-query learning curve for Llama 3.1 8B.
Values are mean $\pm$ sample standard deviation over seeds 42, 2026 and 3407
in percentage points. All checkpoints rerank the same 597 LitSearch queries at
$K_{\mathrm{eval}}=100$.}
\label{tab:parent-budget-detailed}
\end{table}

Tables~\ref{tab:litsearch-subsets-appendix} and
\ref{tab:csfcube-facet-gains-appendix} summarize the subset and facet results,
while Table~\ref{tab:litsearch-ablation-detailed} reports the complete LitSearch
ablation breakdown by first-stage retriever.

\begin{table}[htbp]
\centering
\scriptsize
\resizebox{\columnwidth}{!}{
\begin{adjustbox}{max width=\linewidth}
\begin{tabular}{l c c c c c c}
\hline
Configuration & IB R@20 & IS R@5 & IS R@20 & AB R@20 & AS R@5 & AS R@20 \\
\hline
Llama 3.1 8B FT + Qwen3  & 69.7 & 67.7 & 75.8 & 82.9 & 80.1 & 88.6 \\
Gemma 3 12B FT + Qwen3   & 70.6 & 68.4 & 78.8 & 82.9 & 82.9 & 89.6 \\
Qwen 2.5 14B FT + Qwen3  & 70.2 & \textbf{69.3} & 78.4 & 80.0 & 82.0 & 88.6 \\
Gemma 3 27B FT + Qwen3   & \textbf{75.2} & \textbf{69.3} & \textbf{81.2} & \textbf{85.7} & \textbf{85.3} & \textbf{91.9} \\
\hline
\end{tabular}
\end{adjustbox}
}
\caption{LitSearch primary cohort results by subset for fine-tuned models with
Qwen3-Embedding-8B as the first-stage retriever. IB/IS denote inline
broad/specific queries, while AB/AS denote author-written broad/specific queries.
Values are percentages.}
\label{tab:litsearch-subsets-appendix}
\end{table}

\begin{table}[htbp]
\centering
\scriptsize
\setlength{\tabcolsep}{3pt}
\begin{adjustbox}{max width=\linewidth}
\begin{tabular}{l c c c c}
\hline
Model & Background & Method & Result & All \\
\hline
Llama 3.1 8B     & +0.063 & +0.036 & \textbf{+0.103} & \textbf{+0.068} \\
% NUM-CHECK [NUM-018] Facet table background maximum is emphasized on the wrong model: Gemma 12B +0.053 bold -> Llama +0.063 bold. If bold marks column maxima, move background emphasis to Llama; also consistently emphasize its Method gain +0.036. The +0.053 value itself is correct.
Gemma 3 12B      & \textbf{+0.053} & +0.013 & +0.031 & +0.032 \\
Qwen 2.5 14B     & +0.047 & +0.007 & -0.020 & +0.010 \\
Gemma 3 27B      & -0.045 & +0.003 & -0.029 & -0.023 \\
\hline
\end{tabular}
\end{adjustbox}
\caption{CSFCube test-split nDCG\%@$20$ gains by facet, computed as fine-tuned minus base. Positive values indicate improvement after fine-tuning.}
\label{tab:csfcube-facet-gains-appendix}
\end{table}

\begin{table}[htbp]
\centering
\scriptsize
\setlength{\tabcolsep}{4pt}
\renewcommand{\arraystretch}{0.98}
\begin{tabularx}{\textwidth}{X l r c c c}
\hline
Comparison & Retriever & $\Delta$R@5 & 95\% CI & $p$ & $p_{\mathrm{H}}$ \\
\hline
\multicolumn{6}{l}{\textit{Base--FT comparisons}} \\
Llama 3.1 8B FT $-$ base & BM25 & 12.34 & [8.99, 15.75] & $<.001$ & .0012 \\
Llama 3.1 8B FT $-$ base & GRIT & 13.32 & [9.60, 16.97] & $<.001$ & .0012 \\
Llama 3.1 8B FT $-$ base & Qwen3 & 16.19 & [12.42, 19.91] & $<.001$ & .0012 \\
Gemma 3 12B FT $-$ base & BM25 & 1.62 & [$-$0.64, 3.88] & .1616 & 1.0000 \\
Gemma 3 12B FT $-$ base & GRIT & 0.47 & [$-$1.95, 2.93] & .7391 & 1.0000 \\
Gemma 3 12B FT $-$ base & Qwen3 & 1.59 & [$-$0.61, 3.88] & .1776 & 1.0000 \\
Qwen 2.5 14B FT $-$ base & BM25 & $-$0.47 & [$-$2.90, 1.84] & .7359 & 1.0000 \\
Qwen 2.5 14B FT $-$ base & GRIT & 0.78 & [$-$1.70, 3.27] & .5495 & 1.0000 \\
Qwen 2.5 14B FT $-$ base & Qwen3 & 1.46 & [$-$0.97, 3.96] & .2549 & 1.0000 \\
Gemma 3 27B FT $-$ base & BM25 & $-$0.11 & [$-$2.71, 2.40] & .9047 & 1.0000 \\
Gemma 3 27B FT $-$ base & GRIT & $-$0.92 & [$-$3.54, 1.69] & .5206 & 1.0000 \\
Gemma 3 27B FT $-$ base & Qwen3 & 1.62 & [$-$1.09, 4.30] & .2479 & 1.0000 \\
\hline
\multicolumn{6}{l}{\textit{Gemma 3 27B FT local-baseline comparisons}} \\
Gemma 3 27B FT $-$ retrieval only & GRIT & 1.73 & [$-$0.39, 3.91] & .1309 & .2618 \\
Gemma 3 27B FT $-$ RankZephyr & GRIT & 1.40 & [$-$1.34, 4.16] & .3212 & .3212 \\
Gemma 3 27B FT $-$ Qwen3-Reranker-8B & GRIT & 10.67 & [7.18, 14.28] & $<.001$ & $<.001$ \\
Gemma 3 27B FT $-$ retrieval only & Qwen3 & 3.49 & [1.48, 5.53] & $<.001$ & .0020 \\
Gemma 3 27B FT $-$ RankZephyr & Qwen3 & 3.07 & [0.42, 5.78] & .0250 & .0750 \\
Gemma 3 27B FT $-$ Qwen3-Reranker-8B & Qwen3 & 12.81 & [9.21, 16.39] & $<.001$ & $<.001$ \\
\hline
\end{tabularx}
\caption{Paired uncertainty for the descriptive LitSearch R@5 endpoint. Deltas
are candidate minus reference in percentage points. Holm correction follows the
same base--FT and local-baseline families used for the primary metrics.}
\label{tab:litsearch-r5-uncertainty}
\end{table}

\begin{table}[htbp]
\centering
\scriptsize
\setlength{\tabcolsep}{3.2pt}
\renewcommand{\arraystretch}{1.06}
\begin{adjustbox}{max width=\linewidth}
\begin{tabular}{l c c c c c c c c c}
\hline
\multirow{2}{*}{Configuration}
 & \multicolumn{3}{c}{BM25}
 & \multicolumn{3}{c}{GRIT}
 & \multicolumn{3}{c}{Qwen3} \\
 & R@20 & R@5 & nDCG@20
 & R@20 & R@5 & nDCG@20
 & R@20 & R@5 & nDCG@20 \\
\hline
Standard: $V=3,K=50,P=5$
 & \textbf{67.65} & 55.37 & 50.52
 & \textbf{82.44} & 70.14 & \textbf{64.99}
 & \textbf{81.93} & 69.30 & 65.11 \\
$V=1,K=50,P=5$
 & 66.65 & 55.45 & \textbf{50.57}
 & 80.85 & \textbf{71.17} & 64.51
 & 81.68 & \textbf{70.55} & 65.15 \\
$V=3,K=20,P=5$
 & 65.56 & \textbf{56.37} & 49.55
 & 81.52 & 70.64 & 64.71
 & 80.89 & 69.94 & \textbf{65.26} \\
$V=3,K=100,P=5$
 & 65.51 & 52.60 & 49.43
 & 78.09 & 67.21 & 62.69
 & 76.65 & 64.86 & 61.81 \\
$V=3,K=50,P=1$
 & 64.14 & 52.02 & 48.88
 & 81.77 & 70.03 & 63.81
 & 81.79 & 70.11 & 64.80 \\
$V=3,K=50,P=10$
 & 65.39 & 52.60 & 48.46
 & 80.82 & 69.66 & 64.16
 & 79.80 & 69.55 & 64.53 \\
No pool augmentation
 & 65.06 & 54.30 & 50.27
 & 80.13 & 69.11 & 64.81
 & 79.78 & 68.43 & 64.97 \\
No aspect decomposition
 & 57.94 & 43.81 & 40.78
 & 80.01 & 69.13 & 62.25
 & 80.50 & 69.02 & 62.97 \\
\hline
\end{tabular}
\end{adjustbox}
\caption{Complete reconstructed-environment LitSearch ablations for Llama 3.1 8B,
separated by first-stage retriever, at $K_{\mathrm{eval}}=100$. All values are
overall averages over the same 597 queries and are reported as percentages.
The no-pool-augmentation control uses one pool without
structured mode cycling. Bold indicates the best value for each retriever and
metric.}
\label{tab:litsearch-ablation-detailed}
\end{table}

\FloatBarrier
\section{Compact-Reranker Paired Comparisons}
\label{app:compact-uncertainty}

Table~\ref{tab:compact-uncertainty} reports both the comparison with the
primary Llama FT checkpoint and the effect of adapting GTE-ModernBERT itself.
LitSearch uses paired resampling over 597 queries. CSFCube uses the pooled
50 cases with resampling clustered over 34 query papers, so its case-weighted
means differ slightly from the released equal-fold results in
the main article, Table~16. All deltas and intervals in this table are
expressed in percentage points, including the CSFCube metrics.

\begin{table}[htbp]
\centering
\scriptsize
\setlength{\tabcolsep}{4pt}
\renewcommand{\arraystretch}{0.98}
\begin{adjustbox}{max width=\linewidth}
\begin{tabular}{l l l r c c c}
\hline
Comparison & Retriever / dataset & Metric & $\Delta$ & 95\% CI & $p$ & $p_{\mathrm{H}}$ \\
\hline
Llama FT $-$ GTE adapted & BM25 & R@20 & $-$5.30 & [$-$7.45, $-$3.24] & $<.001$ & $<.001$ \\
Llama FT $-$ GTE adapted & BM25 & nDCG@20 & $-$10.73 & [$-$12.89, $-$8.67] & $<.001$ & $<.001$ \\
Llama FT $-$ GTE adapted & BM25 & R@5 & $-$10.75 & [$-$13.82, $-$7.68] & $<.001$ & $<.001$ \\
Llama FT $-$ GTE adapted & GRIT & R@20 & $-$3.14 & [$-$5.58, $-$0.68] & .0176 & .0176 \\
Llama FT $-$ GTE adapted & GRIT & nDCG@20 & $-$2.09 & [$-$4.40, 0.08] & .0625 & .1250 \\
Llama FT $-$ GTE adapted & GRIT & R@5 & $-$0.94 & [$-$4.08, 2.16] & .5688 & 1.0000 \\
Llama FT $-$ GTE adapted & Qwen3 & R@20 & $-$3.48 & [$-$5.92, $-$1.11] & .0052 & .0104 \\
Llama FT $-$ GTE adapted & Qwen3 & nDCG@20 & $-$1.67 & [$-$3.97, 0.59] & .1619 & .1619 \\
Llama FT $-$ GTE adapted & Qwen3 & R@5 & $-$0.73 & [$-$3.77, 2.32] & .6442 & 1.0000 \\
Llama FT $-$ GTE adapted & CSFCube (pooled) & R@20 & $-$23.90 & [$-$30.99, $-$17.06] & $<.001$ & $<.001$ \\
Llama FT $-$ GTE adapted & CSFCube (pooled) & nDCG\%@20 & $-$24.19 & [$-$30.29, $-$18.04] & $<.001$ & $<.001$ \\
\hline
GTE adapted $-$ zero-shot & BM25 & R@20 & $-$0.17 & [$-$0.84, 0.34] & 1.0000 & 1.0000 \\
GTE adapted $-$ zero-shot & BM25 & nDCG@20 & $-$0.06 & [$-$0.70, 0.59] & .8537 & .9667 \\
GTE adapted $-$ zero-shot & BM25 & R@5 & $-$0.73 & [$-$2.01, 0.50] & .2994 & .4976 \\
GTE adapted $-$ zero-shot & GRIT & R@20 & 0.28 & [$-$0.87, 1.45] & .6709 & 1.0000 \\
GTE adapted $-$ zero-shot & GRIT & nDCG@20 & 0.30 & [$-$0.53, 1.11] & .4834 & .9667 \\
GTE adapted $-$ zero-shot & GRIT & R@5 & 1.23 & [$-$0.22, 2.74] & .1179 & .3537 \\
GTE adapted $-$ zero-shot & Qwen3 & R@20 & 1.20 & [0.00, 2.46] & .0592 & .1776 \\
GTE adapted $-$ zero-shot & Qwen3 & nDCG@20 & 0.79 & [$-$0.01, 1.60] & .0569 & .1707 \\
GTE adapted $-$ zero-shot & Qwen3 & R@5 & 0.98 & [$-$0.67, 2.60] & .2488 & .4976 \\
GTE adapted $-$ zero-shot & CSFCube (pooled) & R@20 & 6.97 & [3.66, 10.60] & $<.001$ & $<.001$ \\
GTE adapted $-$ zero-shot & CSFCube (pooled) & nDCG\%@20 & 2.97 & [1.58, 4.38] & $<.001$ & $<.001$ \\
\hline

\end{tabular}

\end{adjustbox}
\caption{Paired uncertainty for the compact reranker. Deltas follow the
subtraction in the first column; negative Llama-minus-GTE values favor GTE.
Intervals are pointwise 95\% bootstrap intervals. Holm correction is applied
separately within the three retriever comparisons for each LitSearch metric and
contrast. Each CSFCube contrast has one comparison per metric. R@5 remains a
descriptive secondary endpoint.}
\label{tab:compact-uncertainty}
\end{table}

\FloatBarrier
\section{Reconstructed-Environment Replication}
\label{app:reconstructed-replication}

After reconstructing the software environment, we replicate all four model
families using 16-bit LoRA. Table~\ref{tab:reconstructed-main-replication} reports
the results. They provide evidence of qualitative reproducibility and connect the
matched ablations with every model family. They do not replace the
original-cohort results and we do not combine the two cohorts in confidence
intervals or significance tests.

\begin{table}[htbp]
\centering
\scriptsize
\setlength{\tabcolsep}{4pt}
\renewcommand{\arraystretch}{1.06}
\begin{adjustbox}{max width=\linewidth}
\begin{tabular}{l l c c c}
\hline
Model & Retriever & R@20 & R@5 & nDCG@20 \\
\hline
\multirow{3}{*}{Llama 3.1 8B FT}
 & BM25  & 67.65 & 55.37 & 50.52 \\
 & GRIT  & 82.44 & 70.14 & 64.99 \\
 & Qwen3 & 81.93 & 69.30 & 65.11 \\
\hline
\multirow{3}{*}{Gemma 3 12B FT}
 & BM25  & 67.91 & 53.94 & 49.70 \\
 & GRIT  & 81.27 & 69.19 & 64.11 \\
 & Qwen3 & 82.07 & 69.27 & 64.31 \\
\hline
\multirow{3}{*}{Qwen 2.5 14B FT}
 & BM25  & 64.00 & 53.08 & 47.67 \\
 & GRIT  & 80.15 & 72.37 & 64.45 \\
 & Qwen3 & 79.58 & 69.77 & 63.94 \\
\hline
\multirow{3}{*}{Gemma 3 27B FT (16-bit)}
 & BM25  & 67.43 & 56.68 & 50.67 \\
 & GRIT  & 82.69 & 72.04 & 64.83 \\
 & Qwen3 & 82.79 & 71.95 & 65.36 \\
\hline
\end{tabular}
\end{adjustbox}
\caption{Reconstructed-environment replication of all four fine-tuned model
families on the 597-query LitSearch evaluation at
$K_{\mathrm{eval}}=100$. We report percentages from one canonical run for each
model. Table~\ref{tab:seed-stability} provides the three-run statistics for Llama
and Gemma 3 12B.}
\label{tab:reconstructed-main-replication}
\end{table}

\FloatBarrier
\section{Exploratory Variants Excluded from the Final Pipeline}

We explored three additional variants but did not retain them in the main training
recipe. All three are relevant to the broader design space. However, they either
failed to improve the final system or produced gains that were too inconsistent for
inclusion in the main pipeline. These experiments use the original-environment
checkpoints, so we do not compare them numerically with the
reconstructed-environment experiments.

\subsection{Alternative efficiency-oriented formulation via single-token decoding}

We also implement an ablation experiment inspired by \textit{Faster improved listwise reranking with single token
decoding} (FIRST)~\cite{Reddy}. Instead of generating a multi-token permutation
autoregressively, FIRST obtains the complete ranking in one forward pass by
reading the last-position logits at the $K$ candidate-label positions.

\paragraph{Identifier design} We assign a short label to each candidate
(\texttt{A}, \texttt{B}, \dots, \texttt{Z}, \texttt{AA}, \dots). Every label
must encode as exactly one token under the target model's tokenizer, which we
verify before training. All labels are single tokens with the Gemma~3 tokenizer,
comfortably covering our $K = 50$ setting.

\paragraph{Training objective} From the last-position logits, we gather the values
at the $K$ label token IDs and obtain the score vector
$z \in \mathbb{R}^K$. We then apply a \emph{gain-weighted ListMLE}
loss~\cite{Xia}:
\[
\mathcal{L}(z, \pi) = \sum_{i=1}^{K} g_{\pi(i)} \left(
\log \sum_{j \ge i} \exp(z_{\pi(j)}) - z_{\pi(i)} \right),
\]
where $\pi$ denotes the ideal relevance-sorted permutation and
$g_r = 2^r - 1$ denotes the DCG-style gain derived from each 0-2 graded label.
Negative candidates ($r = 0$) contribute no gain, but remain in the partition
function. We normalize each sample loss by its total gain
$\sum_i g_{\pi(i)}$ and average the mini-batch loss over samples.

\paragraph{Training setup} We train this variant for $2$ epochs at
$2 \times 10^{-5}$, with a per-device batch size of $1$ and gradient
accumulation of $16$ (effective batch size $16$).

\paragraph{Inference} At inference time, we run the model once on the reranking
prompt and gather the last-position logits for the $K$ label token IDs. A
descending \texttt{argsort} over this $K$-dimensional vector gives the ranking.
The approach reduces reranking from $O(K)$ decoded tokens to one forward pass.

\begin{table}[htbp]
\centering
\renewcommand{\arraystretch}{1.1}
\setlength{\tabcolsep}{4pt}
\footnotesize
\begin{adjustbox}{max width=\linewidth}
\begin{tabular}{l c c c}
\hline
Configuration & R@20 & R@5 & nDCG@20 \\
\hline
Gemma 3 12B FT, generative (GRIT) & \textbf{81.2} & \textbf{70.2} & \textbf{64.1} \\
Gemma 3 12B FT, FIRST (GRIT)      & 70.8 & 35.8 & 35.0 \\
\hline
\end{tabular}
\end{adjustbox}
\caption{LitSearch overall-average comparison of our final generative listwise
reranker and the FIRST-inspired single-token decoding variant. Both use Gemma 3 12B with
GRIT as the first-stage retriever. Values are percentages.}
\label{tab:first-ablation}
\end{table}

As Table~\ref{tab:first-ablation} shows, FIRST performs substantially worse than
the final generative listwise method, particularly on top-rank metrics. Its single
forward pass makes inference more efficient, but the loss in ranking quality is too
large for our scientific reranking setting. We therefore exclude FIRST from the
main comparisons and the final pipeline.

\subsection{Prompt repetition at inference time}

\begin{table}[htbp]
\centering
\renewcommand{\arraystretch}{1.1}
\setlength{\tabcolsep}{3.5pt}
\footnotesize
\begin{adjustbox}{max width=\linewidth}
\begin{tabular}{l l c c c c}
\hline
Model & Retriever & R@20 & R@5 & nDCG@20 & nDCG@5 \\
\hline
\multicolumn{6}{l}{\textit{Llama 3.1 8B}} \\
Base & BM25  & 57.86 & 40.32 & 40.69 & 35.18 \\
FT   & BM25  & \textbf{65.56} & \textbf{52.80} & \textbf{48.38} & \textbf{44.44} \\
Base & GRIT  & 79.18 & 55.57 & 56.73 & 48.96 \\
FT   & GRIT  & \textbf{80.13} & \textbf{67.77} & \textbf{63.10} & \textbf{59.28} \\
Base & Qwen3 & \textbf{80.03} & 56.90 & 56.28 & 48.70 \\
FT   & Qwen3 & 79.69 & \textbf{67.68} & \textbf{63.49} & \textbf{59.79} \\
\hline
\multicolumn{6}{l}{\textit{Gemma 3 27B}} \\
Base & GRIT  & 83.23 & 73.49 & 67.65 & 64.73 \\
Base & Qwen3 & 80.94 & 69.60 & 64.48 & 61.08 \\
FT   & Qwen3 & \textbf{82.40} & \textbf{71.14} & \textbf{65.88} & \textbf{62.49} \\
\hline
\end{tabular}
\end{adjustbox}
\caption{Preliminary LitSearch overall-average results with prompt repetition at
inference time. Values are percentages over all 597 queries. We report these runs
only as exploratory evidence and do not include them in the main comparison
tables.}
\label{tab:prompt-repetition}
\end{table}

Following Leviathan et al.~\cite{Leviathan}, we also study prompt repetition as
an inference-only variant. Their work shows that repeating the input can improve
accuracy for non-reasoning LLM calls without changing the output format. In our
setting, we duplicate the complete RankGPT-style reranking prompt before
generation:
\[
p(q, D) \rightarrow p(q, D)\;p(q, D),
\]
where $p(q, D)$ is the prompt containing query $q$ and candidate list $D$. The
instruction still requests only a permutation of bracketed candidate identifiers,
so we can use the same output parser.

Prompt repetition requires no additional training and is model-agnostic. These are
useful properties. However, it also roughly doubles the input length during
prefill, which is costly for listwise prompts that can already contain up to 100
title-abstract candidates. The method is therefore less suitable for full-paper
reranking or for inputs near the model's context limit.

In preliminary LitSearch runs, prompt repetition improved several top-rank
metrics, particularly for zero-shot base rerankers. The effect is not uniformly
positive across model-retriever combinations. We therefore exclude repetition
from the final pipeline and treat it as a possible inference-time robustness
technique. A more controlled study should measure latency, context length and
output-parse failures.



\FloatBarrier

\section{Parent-grouped hyperparameter selection details}
\label{supp:grouped-selection-details}

Table~\ref{tab:grouped-hpo-selection} reports the validation results for the
parent-grouped learning-rate search. The rate $5\times10^{-6}$ gives the
lowest parent-averaged NLL on each of the three splits and is selected for the
final refit. The grouped search therefore does not retain the historical
reference rate $1\times10^{-5}$. We adopt its selected rate for the additional
refit. Only the learning rate is varied in this study.
The standard deviations describe variation across splits whose training sets
overlap. They are not standard errors over independent experiments.

\begin{table}[t]
\centering
\scriptsize
\setlength{\tabcolsep}{4pt}
\begin{adjustbox}{max width=\linewidth}
\begin{tabular}{lrr}
\hline
Learning rate & Mean validation NLL & SD across splits \\
\hline
$5\times10^{-6}$ & $\mathbf{0.070711}$ & 0.009043 \\
$1\times10^{-5}$ & 0.072112 & 0.009905 \\
$2\times10^{-5}$ & 0.074852 & 0.008772 \\
\hline
\end{tabular}
\end{adjustbox}
\caption{Parent-query-disjoint learning-rate selection for Llama 3.1 8B over three split assignments. Lower parent-averaged validation NLL is better}
\label{tab:grouped-hpo-selection}
\end{table}

After selection, the chosen checkpoints obtain internal-test parent-averaged
NLLs of 0.076731, 0.076040 and 0.076456 for split seeds 42, 2026 and 3407,
respectively. These evaluations share the same ten held-out parents and do not
constitute independent test samples. The final checkpoint is then trained from
the base model on all 4,494 lists for two epochs using the selected rate.

\FloatBarrier
\section{Tie-handling sensitivity details}
\label{supp:tie-details}

To test whether arbitrary ordering within relevance ties affects adaptation, we
train Llama 3.1 8B with two target policies. The deterministic policy orders
equal-grade candidates by their position in the presented input list. The
epoch-randomized policy independently shuffles candidates within each relevance
grade at each logical epoch, while retaining descending order between grades.
Both policies use response-token cross-entropy. Thus, the experiment changes
the ordering of tied targets, not the relevance labels or the training loss.

We use the same frozen 4,494-list artifact as the compact-reranker study in
Section~\ref{sec:compact-training}, retaining its candidate lists. For
each seed in $\{42,2026,3407\}$, the two policies share candidate membership,
input order, example schedule, base checkpoint and optimization settings.
The runs use unquantized BF16 LoRA with rank 128, alpha 256, dropout 0.05,
learning rate $10^{-5}$, effective batch size two and maximum sequence length
16,384. Two logical epochs give 4,494 optimizer steps. These are newly trained
matched controls. Exact training-list identity with the historical LLM cohort
is not assumed.

Every checkpoint reranks the same 597 LitSearch queries with BM25, GRIT and
Qwen3 candidates at $K_{\mathrm{eval}}=100$. Generation uses temperature zero,
top-$p=1$, seed zero and a 1,024-token output budget. Quality metrics use the
same first-occurrence and missing-ID completion rule described in
Section~\ref{supp:output-parsing}.
We measure validity separately on the raw responses.

Table~8 in the main article shows that randomizing ties does not improve
ranking quality consistently. BM25 R@5 increases descriptively, but nDCG@20
decreases with all three retrievers. For paired inference, we first average
each query's randomized-minus-deterministic difference over the three seeds.
For the retriever average, we also average over BM25, GRIT and Qwen3 before
resampling queries. We use 10,000 paired query-bootstrap replicates and
10,000 sign-flip permutations. Holm correction forms a separate family for
each metric over the three retriever-specific contrasts and their average.
These intervals condition on the retained checkpoints. The sample standard
deviations in Table~8 in the main article describe training-seed variation
separately.

\begin{table}[ht]
\centering
\scriptsize
\setlength{\tabcolsep}{3pt}
\renewcommand{\arraystretch}{1.08}
\begin{adjustbox}{max width=\linewidth}
\begin{tabular}{l r c c}
\hline
Metric & $\Delta$ & 95\% CI & $p_{\mathrm{H}}$ \\
\hline
R@20 & $-$1.21 & [$-$2.41, $-$0.04] & .1404 \\
R@5 & $-$4.81 & [$-$6.67, $-$3.03] & .0004 \\
nDCG@20 & $-$8.40 & [$-$9.91, $-$6.90] & .0004 \\
\hline
\end{tabular}
\end{adjustbox}
\caption{Epoch-randomized minus deterministic tie-policy differences, averaged over retrievers and seeds, in percentage points. Intervals are pointwise; $p_{\mathrm{H}}$ applies within-metric Holm correction}
\label{tab:tie-sensitivity-uncertainty}
\end{table}

The aggregate R@5 and nDCG@20 declines remain significant after correction
(Table~\ref{tab:tie-sensitivity-uncertainty}). The R@20 contrast does not,
despite its unadjusted confidence interval excluding zero. This is not evidence
of equivalence between the policies.

Output validity changes substantially. Averaging over the equally sized
retriever--seed groups, the raw exact-permutation rate falls from 52.67\% with
deterministic ties to 0.24\% with randomized ties. The mean number of missing
candidate IDs rises from 1.42 to 32.25. Both conditions have zero complete
fallback, but this only means that some valid IDs were emitted. It does not
imply complete raw rankings. These differences qualify the interpretation of
the repaired quality scores.

We also audit whether the output budget directly limits the measured top-rank
metrics. Of 1,349 length-finished responses across the six runs, 1,348 already
contain at least 20 distinct valid IDs. In the remaining response, the sole
relevant document is already ranked first. Under the first-occurrence parser,
extending these same prefixes cannot change R@5, R@20 or nDCG@20. This is an
audit of the saved responses, not a rerun with a larger generation budget.

\FloatBarrier
\section{GPT-4o pipeline comparison visualization}
\label{supp:gpt4o-figure}

Figure~\ref{fig:gpt4o-comparison} visualizes selected pipelines from Table~10
in the main article. The pipelines use different first-stage retrievers.
GPT-4o was not rerun, and its per-query outputs are unavailable; this is a
descriptive comparison without paired significance tests.

\begin{figure}[htbp]
  \centering
  \manuscriptfigure[width=\columnwidth]{fig_gpt4o_comparison.pdf}
  \caption{Descriptive comparison of the GPT-4o + GritLM pipeline reported by
  LitSearch~\cite{Ajith}, Qwen3 retrieval without reranking, and Qwen3 with
  fine-tuned Gemma 3 27B.}
  \label{fig:gpt4o-comparison}
\end{figure}

\FloatBarrier

\begin{thebibliography}{99}

\bibitem{Ajith}Ajith, A., Xia, M., Chevalier, A., Goyal, T., Chen, D., \& Gao, T. (2024). LitSearch: A retrieval benchmark for scientific literature search. In Proceedings of the 2024 Conference on Empirical Methods in Natural Language Processing (pp. 15068--15083). doi: \url{https://doi.org/10.18653/v1/2024.emnlp-main.840}.

\bibitem{Mysore}Mysore, S., O'Gorman, T., McCallum, A., \& Zamani, H. (2021). CSFCube---A test collection of computer science research articles for faceted query by example. arXiv preprint arXiv:2103.12906. doi: \url{https://doi.org/10.48550/arXiv.2103.12906}.

\bibitem{Sun}Sun, W., Yan, L., Ma, X., Wang, S., Ren, P., Chen, Z., et al. (2023). Is ChatGPT good at search? Investigating large language models as re-ranking agents. In Proceedings of the 2023 Conference on Empirical Methods in Natural Language Processing (pp. 14918--14937). doi: \url{https://doi.org/10.18653/v1/2023.emnlp-main.923}.

\bibitem{Xia}Xia, F., Liu, T. Y., Wang, J., Zhang, W., \& Li, H. (2008). Listwise approach to learning to rank: Theory and algorithm. In Proceedings of the 25th International Conference on Machine Learning (pp. 1192--1199). doi: \url{https://doi.org/10.1145/1390156.1390306}.

\bibitem{Reddy}Gangi Reddy, R., Doo, J., Xu, Y., Sultan, M. A., Swain, D., Sil, A., et al. (2024). FIRST: Faster improved listwise reranking with single token decoding. In Proceedings of the 2024 Conference on Empirical Methods in Natural Language Processing (pp. 8642--8652). doi: \url{https://doi.org/10.18653/v1/2024.emnlp-main.491}.

\bibitem{Zhang}Zhang, Y., Li, M., Long, D., Zhang, X., Lin, H., Yang, B., et al. (2025). Qwen3 Embedding: Advancing text embedding and reranking through foundation models. arXiv preprint arXiv:2506.05176. doi: \url{https://doi.org/10.48550/arXiv.2506.05176}.

\bibitem{Hu}Hu, E. J., Shen, Y., Wallis, P., Allen-Zhu, Z., Li, Y., Wang, S., et al. (2022). LoRA: Low-rank adaptation of large language models. In International Conference on Learning Representations (ICLR 2022). arXiv version doi: \url{https://doi.org/10.48550/arXiv.2106.09685}.

\bibitem{Dettmers}Dettmers, T., Pagnoni, A., Holtzman, A., \& Zettlemoyer, L. (2023). QLoRA: Efficient finetuning of quantized LLMs. Advances in Neural Information Processing Systems, 36, 10088--10115. doi: \url{https://doi.org/10.52202/075280-0441}.

\bibitem{unsloth}Han, D., Han, M., \& Unsloth Team. (2023). Unsloth [Software]. \url{https://github.com/unslothai/unsloth}.

\bibitem{Kwon}Kwon, W., Li, Z., Zhuang, S., Sheng, Y., Zheng, L., Yu, C. H., et al. (2023). Efficient memory management for large language model serving with PagedAttention. In Proceedings of the 29th Symposium on Operating Systems Principles (pp. 611--626). doi: \url{https://doi.org/10.1145/3600006.3613165}.

\bibitem{Leviathan}Leviathan, Y., Kalman, M., \& Matias, Y. (2025). Prompt repetition improves non-reasoning LLMs. arXiv preprint arXiv:2512.14982. doi: \url{https://doi.org/10.48550/arXiv.2512.14982}.

\bibitem{gte-modernbert}Tongyi Lab, Alibaba Group. (n.d.). GTE-reranker-modernbert-base [Model card]. \url{https://huggingface.co/Alibaba-NLP/gte-reranker-modernbert-base}. Evaluated revision: \texttt{204bdd4}.

\end{thebibliography}

\end{document}
